\section{Configuration reference}
\subsection{Current cross-talk entry-script settings}
These values describe the inspected script, not a recommended preset for an unverified setup. Position, range and offset are sent directly to \code{CHANnel2:POSition}, \code{CHANnel2:RANGe} and \code{CHANnel2:OFFset}. Consult the scope command reference before translating a front-panel V/div setting into a range value.
\begin{center}
\begin{tabular}{@{}lrrr@{}}
\toprule
Pair category & Position & Range & Offset \\
\midrule
Diagonal ($i=j$) & 0 & 5 & 0 \\
Adjacent numbers ($|i-j|=1$) & 0 & 0.3 & 0 \\
Other pairs & 0 & 0.1 & 0 \\
\bottomrule
\end{tabular}
\end{center}
The shared defaults differ: diagonal $(-4,5,-1.8)$, adjacent $(0,0.15,0.009)$ and other $(0,0.08,0.009)$, in position/range/offset order. The entry script overrides all nine values.

\subsection{Tower maps}
\label{sec:tower-maps}
For towers 2, 3, 5, 7, 9, 11, 12 and 14, the current display order is
\begin{lstlisting}
[9, 13, 2, 5, 11, 14, 3, 8, 6, 12, 4, 10, 7, 1, 15]
\end{lstlisting}
and the study-cell-to-channel map is
\begin{lstlisting}
[9, 13, 2, 5, 11, 14, 3, 6, 12, 4, 10, 7, 1, 15, 8]
\end{lstlisting}
For towers 4, 6, 8, 10 and 13, display order is \code{1:15}, and the study map is \code{[1:7, 9:15, 8]}. T0 uses \code{1:15} for both maps. Other tower numbers produce an error until their mappings are added.

\subsection{Multiplexer command map}
Before selecting a pair, the acquisition routines send \code{C}. Injection channels 1--7 use \code{U 8 3} followed by \code{L 0 <channel>}; channels 8--15 use \code{L 8 3} followed by \code{L 1 <channel-8>}. Readout channels 1--7 use \code{U 8 5} followed by \code{L 6 <channel>}; channels 8--15 use \code{L 8 5} followed by \code{L 7 <channel-8>}.

These commands record the implemented routing, not an independently verified connector pinout. Check \code{set_matrix_channel.m} against the real wiring after any hardware change.

\subsection{Scope waveform transfer}
\code{read_scope.m} requests \code{CHAN2:DATA:HEAD?}, then selects \code{:FORMAT ASCii} and requests \code{CHAN2:DATA?}. It interprets the first three header fields as start time, end time and sample count, constructs a linearly spaced time vector and checks the returned sample count. This command dialect and header interpretation must be checked when replacing the oscilloscope.

\section{Measurement record template}
\label{sec:record}
Keep a completed record alongside each campaign's raw data. The following is a proposed template, not an automatically generated toolkit output.
\begin{longtable}{@{}>{\raggedright\arraybackslash}p{0.33\linewidth}>{\raggedright\arraybackslash}p{0.62\linewidth}@{}}
\toprule
Field & Value to enter \\
\midrule
\endhead
Date and operator & Date, local time/time zone, operator name. \\
PCB and tower & Board identifier/revision, tower number, connector orientation, photograph reference. \\
Measurement & Cross-talk, T0 capacitance or tower capacitance. \\
Software & Toolkit revision, saved entry-script copy, MATLAB/toolbox versions, Windows and driver versions. \\
Instruments & Model, serial number, identity response and VISA resource for matrix and scope; generator identifier. \\
Generator & Waveform, frequency, amplitude, offset, output impedance/load convention and cabling. \\
Scope & Input impedance/coupling, probe factor, bandwidth, trigger, timebase, sample rate, record length and CH1 settings. \\
CH2 presets & Position, range and offset for every applicable pair category. \\
Timing & Pre-acquisition settle time and post-acquisition read delay. \\
Reference & T0 dataset location, acquisition time, resistance value and measurement source. \\
Data and analysis & Data root, dataset label, overwrite/resume choice, output label, shaping times and geometry version. \\
Observations & Diagnostic waveform paths, warnings, rejected/repeated channels, setup changes and reasons. \\
Archive & Raw-data location, configuration record and exported-result location. \\
\bottomrule
\end{longtable}

\section{Source notes and remaining work}
\subsection{Local sources used for this draft}
\begin{itemize}
\item Historical installation note: \path{/Users/huang/workarea/Documentation/FCC/Install_GPIB_Windows.md}.
\item Toolkit: \path{/Users/huang/cernbox/Allegro_PCB/pcb_test_toolkit}, including its README, entry scripts, configuration, acquisition, waveform-loading, shaping, fitting and analysis functions.
\item Checkout HEAD at drafting: \code{ccf13d3ce807346d0af59ee9a0af418b9924dc9a}. The working tree had existing modifications in \code{analyze_capacity.m}, \code{default_config.m}, \code{fit_tau.m} and \code{README.md}. Descriptions refer to these inspected working files, so this commit alone does not fully identify the documented software snapshot.
\end{itemize}
The original toolkit and historical documentation were not modified. This document does not replace instrument manuals or a validated bench record.

\subsection{Priority additions from practical experience}
\begin{enumerate}
\item An annotated wiring diagram, instrument inventory, connector photographs and T0 description.
\item Verified generator and scope settings for each measurement type.
\item Representative good and bad waveforms with explanations and acceptance criteria.
\item Physical channel maps, the reason for omitting position/channel 8, and geometry sources.
\item Reference-repeatability checks, resistance uncertainty and fit-quality criteria.
\item Known bench problems with confirmed remedies, and the long-term data/software archive location.
\end{enumerate}
